\section{Introduction}

Fungal mitochondrial genomes are unusually plastic. Within a single genus they routinely vary
several-fold in length while encoding essentially the same small set of proteins, and the
standard explanation is mobile intron content: group I introns, and the homing endonucleases
they carry, invade conserved genes and expand the genome without adding respiratory function
\cite{zaccaron2020bipolaris}.

That explanation is widely repeated and less often tested within a genus with a consistent
annotation. The difficulty is that deposited mitogenome annotations differ enormously in
completeness --- the companion paper to this one shows that the reference used here carried
thirteen features where our pipeline recovers fifty-three --- so comparing published feature
counts measures annotation effort as much as biology.

We therefore annotated twelve complete \textit{Pleurotus} mitogenomes with a single pipeline
and asked: does intron content explain size; is any of the sequence duplicated; can the
uncharacterised open reading frames be assigned a function or at least shown to be real; and,
since mitochondrial function is jointly encoded by two genomes, is nuclear-encoded
mitochondrial expression coordinated with the mitochondrial genome's own output across tissues.

\section{Results}

\subsection{Mitogenome content}

The \textit{P.\ ostreatus} BOM\_ss5 mitochondrion is 71,949\,bp at 26.6\% GC and, under the
corrected annotation, carries 20 protein-coding genes, 25 tRNAs, 2 rRNAs and 6 spliced introns
(Fig.~\ref{fig:circ}). Coverage is markedly uneven: the ribosomal RNA regions reach $10^6$
while much of the protein-coding complement sits between $10^1$ and $10^3$.

\subsection{Intron content tracks genome size, within and between species}

Across twelve mitogenomes, length spans 60,694\,bp (\textit{P.\ citrinopileatus}) to
102,950\,bp (\textit{P.\ giganteus}), a 1.7-fold range with no corresponding change in the
respiratory gene set (Table~\ref{tab:comparative}).

Detected group I intron content correlates with length at $r=+0.89$, and ORF count more
strongly at $r=+0.96$ (Fig.~\ref{fig:comp}). The second correlation is not independent of the
first: group I introns encode homing endonucleases, so an intron gain adds both sequence and an
ORF. Together they support the standard account.

The effect operates below the species level. Among the seven \textit{P.\ ostreatus} mitogenomes
here, length ranges from 65,277\,bp (PC9.15) to 73,085\,bp (N009) --- a 7.8\,kb spread within
one species, with the smaller genomes carrying three to four detected group I intron hits and
the larger seven to eight. Our own strain sits at the larger end.

Detected intron span is a lower bound: covariance models capture the conserved core rather than
the whole intron. For BOM\_ss5 the models detect 1,123\,bp where split-read evidence in the
companion paper puts the true spliced total at 4,332\,bp, roughly fourfold higher. Detected
span is therefore a comparable index across genomes, not an absolute measurement, and it does
not by itself account for the whole size difference.

\subsection{Little duplication, but what there is duplicates an unknown gene}

Self-alignment recovers only two internal repeats of at least 200\,bp, together 802\,bp or
1.1\% of the genome. This mitogenome is not built from repeated sequence.

Both repeats are nonetheless interesting. All-versus-all comparison of the 51 predicted ORFs
finds four paralogous families. Three are expected: the NADH dehydrogenase subunits nad2, nad4
and nad5 share ancestry rather than representing duplication, and two families of intron-encoded
endonucleases (aI3, aI5, aI8) are homing endonucleases, which proliferate by their nature.

The fourth family is not expected. \texttt{CM148777.1\_26} (64,085--64,510) and
\texttt{CM148777.1\_28} (66,884--67,339) have no Swiss-Prot homologue, and they correspond
exactly to the larger internal repeat at 84.5\% nucleotide identity --- a duplicated
uncharacterised ORF. Both copies are present in seven of eleven other \textit{Pleurotus}
mitogenomes, so the duplication predates the strain, but both are poorly expressed here (five
reads and zero reads respectively), which is equally consistent with a pseudogenised pair.

\subsection{Sensitive search does not name the uncharacterised ORFs, but conservation and
expression identify four candidates}

Twenty predicted ORFs have no Swiss-Prot hit. Scanning them against Pfam --- a profile method
substantially more sensitive than BLAST --- assigns an informative domain to \textbf{none}. The
single hit obtained is to a \textit{Caenorhabditis} neuropeptide-like family at $E=8.8\times
10^{-4}$, which we regard as a low-complexity artefact rather than a functional assignment.

Homology search is therefore exhausted, and we fall back on two orthogonal criteria: is the ORF
conserved in other mitogenomes, and is it transcribed? Four ORFs satisfy both, being present in
at least six of eleven other \textit{Pleurotus} mitogenomes and detected in at least eight of
sixteen libraries (Table~\ref{tab:orfcons}, Fig.~\ref{fig:orf}a). One further ORF is present in
a single other genome and is most likely strain-specific or spurious.

These four are candidate mitochondrial genes that no current database can name. We make no
claim about what they do.

\subsection{A predicted nuclear-encoded mitochondrial proteome}

Reciprocal best hits between the 12,521 nuclear proteins and MitoCarta3.0 give 399 predicted
mitochondrial proteins; a one-way screen would have returned 844, which is why reciprocity was
required. Curated Swiss-Prot mitochondrial localisations independently identify 467. The union
is \textbf{631 proteins, 5.0\% of the proteome}.

The two approaches are complementary by construction. MitoCarta is a mammalian inventory, so
fungal-specific mitochondrial proteins are invisible to it; Swiss-Prot contains curated yeast
and filamentous-fungal entries but covers only 41\% of this proteome at all. The union is
conservative and should be read as a floor.

\subsection{Mito-nuclear balance differs between tissues, but the ranges overlap}

Because mitochondrial function is encoded by two genomes, their relative output is measurable.
Expressing both on the same DESeq2 size factors, the ratio of mtDNA-encoded to nuclear-encoded
mitochondrial expression is highest in the exudophore (median 1.51) and similar across the
other three tissues (0.53--0.74; Table~\ref{tab:retro}, Fig.~\ref{fig:orf}b).

Both exudophore replicates exceed every fuzzy mycelium replicate. They do not exceed every
library: one exuding mycelium library reaches 1.02 and one nodule library 0.94, against 1.14
for the lower exudophore replicate. With $n=2$ for the exudophore and overlapping ranges, this
is an observation to follow up rather than a result.

We flag one methodological point behind it. An earlier version of this comparison used raw
mitochondrial counts against normalised nuclear counts, and reported the highest ratio in
nodule with a correlation of $+0.92$ across tissues. Both were artefacts of the deepest library
sitting in that group. Applying matched size factors moved the maximum to the exudophore and
dropped the correlation to $+0.43$. Ratios between two genomes are only interpretable when both
carry the same normalisation.

\section{Discussion}

The comparative result is a clean confirmation of an accepted model, obtained under a
consistent annotation: within \textit{Pleurotus}, mitogenome size tracks group I intron content
and the ORF load that accompanies it, and the relationship holds within \textit{P.\ ostreatus}
as well as between species. Whether the intron differences between strains reflect recent
mobility or long-standing lineage differences is not answerable from sequence alone.

The more suggestive observations are negative or preliminary. That Pfam names none of twenty
uncharacterised ORFs, while four are both conserved and expressed, is a statement about
databases rather than about the organism: mitochondrial ORFs of this kind are numerous in
fungal mitogenomes \cite{zaccaron2020bipolaris} and remain almost entirely uncharacterised. A
duplicated uncharacterised ORF conserved across seven genomes but essentially untranscribed is
the kind of feature that is easy to annotate away and hard to interpret.

The mito-nuclear ratio is worth pursuing precisely because the exudophore is where the other
analyses in this series also point --- an oxidative, respiratorily active structure. A shifted
balance between the two genomes would be consistent with that. Establishing it needs more
replicates than these libraries provide.

\subsection{Limitations}

Annotation is computational throughout; no gene model has been verified experimentally. The
comparative set is twelve genomes, adequate for a correlation but not for phylogenetically
controlled inference, and the species are related enough that the points are not independent.
Detected intron span underestimates true intron content roughly fourfold and does so unevenly.
The predicted nuclear mito-proteome inherits MitoCarta's mammalian bias and Swiss-Prot's
coverage limits. All tissue statements rest on twelve libraries with the exudophore at $n=2$.
Nothing here demonstrates retrograde signalling; the data speak to expression balance only.

\section{Methods}

Twelve complete \textit{Pleurotus} mitogenomes were retrieved from NCBI and annotated with the
pipeline described in the companion paper: tRNAscan-SE \cite{chan2021trnascan}, barrnap
\cite{seemann2013barrnap}, EMBOSS \texttt{getorf} under genetic code 4, and DIAMOND
\cite{buchfink2015diamond} identification against Swiss-Prot \cite{uniprot2023}. Group I and II
introns were detected with Infernal \cite{nawrocki2013infernal} \texttt{cmsearch} against Rfam
\cite{kalvari2021rfam} models at $E<0.01$. BOM\_ss14, the sibling nucleus of the same dikaryon,
was included as a control and returned counts identical to BOM\_ss5.

Internal repeats and pairwise synteny used \texttt{nucmer} \cite{marcais2018mummer4}
(\texttt{--maxmatch --nosimplify} for self-alignment), filtering the trivial diagonal and
retaining alignments of at least 200\,bp. ORF paralogy used all-versus-all DIAMOND with
single-linkage clustering. Pfam \cite{mistry2021pfam} domains were assigned with
\texttt{hmmscan} \cite{eddy2011hmmer} at $E<10^{-3}$.

MitoCarta3.0 \cite{rath2021mitocarta} human sequences were searched reciprocally against the
nuclear proteome, retaining reciprocal best hits. Swiss-Prot subcellular-location annotations
were parsed from the flatfile. Mitochondrial and nuclear expression were placed on common
DESeq2 \cite{love2014deseq2} size factors before any ratio was computed.

\clearpage
\section*{Figures}

\begin{figure}[htbp]\centering\includegraphics[width=0.82\textwidth]{fig11_circular_map.pdf}
\caption{\textbf{The blue oyster mitochondrial genome.} From the outside: pooled RNA-seq depth
(log scale); protein-coding, ribosomal and transfer RNA features on the forward then reverse
strand; spliced introns; coordinate axis. Interior links join the two internal repeats;
the larger duplicates an uncharacterised ORF at 84.5\% identity.}\label{fig:circ}\end{figure}

\begin{figure}[htbp]\centering\includegraphics[width=\textwidth]{fig12_comparative_size.pdf}
\caption{\textbf{Intron content tracks mitogenome size across \textit{Pleurotus}.}
\textbf{a}, Genome size against detected group I intron span. \textbf{b}, Against ORF count.
All twelve genomes annotated with one pipeline.}\label{fig:comp}\end{figure}

\begin{figure}[htbp]\centering\includegraphics[width=\textwidth]{fig13_orfs_retrograde.pdf}
\caption{\textbf{Uncharacterised ORFs, and mito-nuclear balance.} \textbf{a}, Each
uncharacterised ORF by conservation across other \textit{Pleurotus} mitogenomes and detection
across libraries; point size is read support. \textbf{b}, Ratio of mtDNA-encoded to predicted
nuclear-encoded mitochondrial expression, per library, both on matched size
factors.}\label{fig:orf}\end{figure}

\clearpage
\section*{Tables}
\begin{table}[htbp]
\caption{Complete \textit{Pleurotus} mitogenomes, all annotated with the pipeline used here rather than with their deposited features.}\label{tab:comparative}
\footnotesize
\begin{tabular}{@{}llrrrrrr@{}}
\toprule
Species / strain & Accession & Length (bp) & GC\% & tRNA & ORFs & gpI hits & Intron bp \\
\midrule
P. citrinopileatus & NC\_036998.1 & 60,694 & 29.0 & 24 & 51 & 2 & 933 \\
P. ostreatus PC9.15 & CM057219.1 & 65,277 & 26.2 & 25 & 47 & 3 & 523 \\
P. ostreatus N001 & PX724300.1 & 65,303 & 26.2 & 25 & 47 & 4 & 636 \\
P. ostreatus & OR030114.1 & 65,311 & 26.2 & 25 & 46 & 4 & 636 \\
P. pulmonarius & NC\_061177.1 & 70,671 & 25.3 & 25 & 49 & 6 & 1,367 \\
P. ostreatus BOM\_ss14 & CM148778.1 & 71,947 & 26.6 & 25 & 51 & 5 & 1,295 \\
P. ostreatus BOM\_ss5 (this study) & CM148777.1 & 71,949 & 26.6 & 25 & 51 & 4 & 1,123 \\
P. ostreatus F515 & PX724301.1 & 72,132 & 26.7 & 25 & 54 & 7 & 1,783 \\
P. cornucopiae & NC\_038091.1 & 72,134 & 26.7 & 25 & 53 & 7 & 1,783 \\
P. ostreatus N009 & PX724302.1 & 73,085 & 26.7 & 25 & 55 & 8 & 1,912 \\
P. platypus & NC\_036999.1 & 73,807 & 26.8 & 25 & 55 & 8 & 1,912 \\
P. giganteus & NC\_062374.1 & 102,950 & 25.0 & 25 & 86 & 12 & 2,994 \\
\bottomrule
\end{tabular}
\begin{tablenotes}[flushleft]\footnotesize\item gpI hits are Rfam group I covariance-model matches; the span they cover is a lower bound, since the models capture the conserved core rather than the full intron.\end{tablenotes}
\end{table}

\begin{table}[htbp]
\caption{Uncharacterised mitochondrial ORFs: conservation across \textit{Pleurotus} against expression.}\label{tab:orfcons}
\footnotesize
\begin{tabular}{@{}lrrrc@{}}
\toprule
ORF & Other genomes & Libraries & Reads & In rRNA block \\
\midrule
CM148777.1\_27 & 10/11 & 12/16 & 584 & no \\
CM148777.1\_14 & 10/11 & 11/16 & 524 & no \\
CM148777.1\_44 & 9/11 & 8/16 & 218 & no \\
CM148777.1\_30 & 9/11 & 5/16 & 100 & no \\
CM148777.1\_21 & 9/11 & 4/16 & 60 & yes \\
CM148777.1\_37 & 9/11 & 3/16 & 29 & yes \\
CM148777.1\_43 & 8/11 & 6/16 & 106 & no \\
CM148777.1\_20 & 7/11 & 5/16 & 104 & yes \\
CM148777.1\_28 & 7/11 & 1/16 & 5 & no \\
CM148777.1\_26 & 7/11 & 0/16 & 0 & no \\
CM148777.1\_33 & 6/11 & 11/16 & 376 & no \\
CM148777.1\_51 & 6/11 & 7/16 & 58 & yes \\
CM148777.1\_11 & 6/11 & 3/16 & 30 & no \\
CM148777.1\_50 & 5/11 & 16/16 & 14,070 & yes \\
\bottomrule
\end{tabular}
\begin{tablenotes}[flushleft]\footnotesize\item No Swiss-Prot hit and no informative Pfam domain. An ORF conserved in most genomes and expressed in most libraries is a candidate gene; one present in a single strain is more likely spurious.\end{tablenotes}
\end{table}

\begin{table}[htbp]
\caption{Ratio of mtDNA-encoded to predicted nuclear-encoded mitochondrial expression, per library.}\label{tab:retro}
\footnotesize
\begin{tabular}{@{}lrrrr@{}}
\toprule
Tissue & $n$ & Min & Median & Max \\
\midrule
Exudophore & 2 & 1.139 & 1.507 & 1.875 \\
Nodule & 3 & 0.309 & 0.740 & 0.940 \\
Fuzzy mycelium & 4 & 0.442 & 0.588 & 0.784 \\
Exuding mycelium & 3 & 0.209 & 0.525 & 1.016 \\
\bottomrule
\end{tabular}
\begin{tablenotes}[flushleft]\footnotesize\item Both terms carry the same DESeq2 size factors. Ranges overlap between tissues and the exudophore has $n=2$; this is descriptive.\end{tablenotes}
\end{table}

